\chapter{Ce que la collaboration a produit}
\label{ch:collaborations}

\questionchapitre{À quels besoins, exprimés par qui, ces développements répondent-ils
effectivement ?}


\section{Le cadre de travail}

Les développements présentés dans ce rapport ont été conçus et réalisés dans leur intégralité au
cours du postdoctorat, dans un dialogue continu avec trois interlocuteurs aux apports distincts.

Nicolas Labroche (LIFAT) a assuré la supervision scientifique et encadré l'orientation des
travaux, à commencer par l'arbitrage exposé au chapitre~\ref{ch:contexte}, qui a fixé la
trajectoire de l'année.

Hugo Breuillard, data scientist en charge du projet côté BRGM, a été le point de contact sur les
questions de chaîne de données et de stratégie de prévision. Les échanges ont porté sur
l'articulation entre l'entrepôt développé ici et les besoins de modélisation du programme, et
sur la cohérence des indicateurs produits avec ceux employés par ailleurs.

Augustin Thomas, hydrogéologue au BRGM, a été l'interlocuteur sur l'ensemble des questions de
domaine. C'est de ces échanges que procèdent les décisions qui ancrent la plateforme dans la
pratique hydrogéologique plutôt que dans une transposition générique de méthodes
d'apprentissage.

Cet apport n'est pas resté au stade de la discussion : il est inscrit dans le code, et
identifiable ligne à ligne. Les configurations par famille d'aquifère (tableau~\ref{tab:familles-aquiferes}) et les quatre
critères d'acceptation STOWA (chapitre~\ref{ch:metier}), la règle d'affichage cartographique
(chapitre~\ref{ch:observatoire}) et l'arbitrage sur l'évapotranspiration
(section~\ref{sec:hargreaves}) en sont les quatre traces les plus nettes. Aucune ne procède
d'une procédure automatique, et la dernière mérite d'être relevée pour ce qu'elle dit de la
confrontation disciplinaire : la variable fournie par la réanalyse était disponible,
correctement nommée et directement utilisable, et seule sa confrontation aux ordres de grandeur
connus pour la France a révélé qu'elle ne désignait pas la grandeur attendue. Aucune
vérification interne au code ne pouvait détecter cette erreur, puisqu'il ne s'agissait pas d'une
erreur de calcul.

\section{Les cas d'usage identifiés avec le BRGM}

Les échanges n'ont pas seulement produit des décisions techniques. Ils ont fait émerger une
liste de besoins concrets, qui a servi à orienter les développements et qui reste disponible
pour la suite. Cette liste est un livrable à elle seule : un socle sans usage identifié n'est qu'une
infrastructure en attente.

\begin{table}[H]
\centering
\caption{Les besoins exprimés par les équipes du BRGM, et l'état de leur prise en charge.}
\label{tab:besoins-brgm}
\begin{tabularx}{\textwidth}{@{}L{3.3cm}Yl@{}}
\toprule
\textbf{Besoin exprimé} & \textbf{Ce qu'il recouvre} & \textbf{État} \\
\midrule
Qualifier les piézomètres & Distinguer automatiquement les ouvrages inertiels, réactifs et
influencés par des prélèvements, pour savoir lesquels sont exploitables & outillé \\
Prévoir les niveaux & Anticiper l'évolution d'une nappe à partir du climat & outillé \\
Expliquer les prévisions & Identifier les facteurs déterminants d'une prévision donnée &
outillé \\
Comparer les approches & Confronter méthodes statistiques, modèles métier et apprentissage
profond sur les mêmes données & possible, non fait \\
Explorer les données & Permettre aux experts de parcourir et de comparer les chroniques sans
écrire de requête & outillé \\
Scénariser la sécheresse & Répondre à des questions du type « quelle pluie faudrait-il pour
sortir de cet état ? », en appui des comités sécheresse & outillé hors interface \\
Étudier les échanges nappe-rivière & Détecter et caractériser les relations entre une nappe et
le cours d'eau voisin, notamment le soutien d'étiage & non traité \\
Caractériser la double porosité & Distinguer, au sein d'une même formation, les écoulements
rapides des écoulements lents, que les méthodes actuelles peinent à séparer & non traité \\
Couvrir les nappes captives & Disposer d'un indicateur de situation pour les nappes qui n'en ont
pas, l'indicateur existant ne valant que pour les nappes libres & non traité \\
\bottomrule
\end{tabularx}
\end{table}

La colonne d'état est délibérément sévère : « outillé » signifie
que les moyens existent, pas qu'un résultat a été produit, et « outillé hors interface » signale
que la méthode s'exécute par l'API ou par un carnet de calcul mais reste inaccessible depuis
l'application, pour la raison exposée au chapitre~\ref{ch:apprentissage}. Et les trois dernières lignes sont
assumées comme telles : ce sont des sujets à part entière, identifiés comme prometteurs au cours
des échanges, mais qui dépassaient le périmètre et le temps disponibles. Ils sont consignés ici
parce qu'ils constituent des points d'entrée pour la suite, et parce qu'ils viennent des experts
du domaine plutôt que d'une projection d'informaticien.

Un mot sur la scénarisation. Les équipes du BRGM travaillant
sur la situation des nappes emploient déjà les indicateurs standardisés que produit désormais
l'entrepôt. Coupler ces indicateurs à des scénarios prospectifs, c'est-à-dire répondre non
seulement « où en est-on » mais « que faudrait-il pour en sortir », est le prolongement naturel
de leur pratique actuelle. C'est cette continuité qui rend le socle mobilisable par eux sans
changement de méthode.

\section{L'articulation avec les travaux de recherche en cours}

Les outils d'analyse présentés en section~\ref{sec:xai} ne sont pas isolés. Des travaux menés au
LIFAT sous la direction de Nicolas Labroche portent sur les explications contrefactuelles et sur
la mesure de leur qualité en tant qu'explications, question ouverte que les métriques usuelles ne
tranchent pas.

La contribution de ce postdoctorat à cette thématique est \textbf{d'ordre technique}, et c'est
délibéré : fournir les données, les modèles, les méthodes comparables et les moyens de
validation sans lesquels cette recherche resterait cantonnée à des jeux d'essai génériques. La
répartition est claire, et il vaut mieux l'écrire que la laisser deviner.

\begin{table}[H]
\centering
\caption{Répartition entre ces travaux et la recherche en cours au LIFAT.}
\label{tab:repartition}
\begin{tabularx}{\textwidth}{@{}lX@{}}
\toprule
\textbf{Relève de ces travaux} & Le socle de données, les modèles à expliquer, l'implémentation
des méthodes, la référence de comparaison, le moyen de validation indépendant, et la
qualification des cas d'étude \\
\textbf{Relève des travaux en cours} & La définition de ce qu'est une bonne explication, les
métriques qui la mesurent, et l'évaluation comparative des approches \\
\bottomrule
\end{tabularx}
\end{table}

C'est aussi ce qui explique qu'aucune conclusion méthodologique ne figure dans ce rapport sur la
qualité des contrefactuels produits. Elle n'aurait pas eu de sens : la question relève d'un
travail en cours, et l'anticiper reviendrait à préjuger de son résultat.

\section{Des besoins distincts, un outil commun}

Les deux communautés visées n'attendent pas la même chose, et l'architecture a dû accommoder les
deux sans se scinder.

\begin{table}[H]
\centering
\caption{Deux communautés, deux attentes, un même socle.}
\label{tab:deux-publics}
\begin{tabularx}{\textwidth}{@{}L{3.6cm}YY@{}}
\toprule
 & \textbf{Hydrogéologues du BRGM} & \textbf{Chercheurs en apprentissage} \\
\midrule
\textbf{Question posée} & Où en est la ressource, et que se passe-t-il si l'on prélève ? &
Quelle méthode prédit le mieux, et pourquoi ? \\
\textbf{Attente principale} & Des indicateurs conformes aux pratiques, une scénarisation
crédible & Un jeu de données consolidé, un protocole reproductible \\
\textbf{Obstacle levé} & L'agrégation manuelle de sources hétérogènes & Le coût d'entrée à
l'expérimentation \\
\textbf{Ce qui a été livré} & Observatoire, indicateurs validés, laboratoire de modèles métier,
scénarios bornés & Entrepôt interrogeable, catalogue d'architectures, traçabilité des
expériences, outillage d'explicabilité \\
\bottomrule
\end{tabularx}
\end{table}

Le point de rencontre est l'entrepôt. Parce que les indicateurs métier et les jeux
d'apprentissage sont produits par la même chaîne et validés une seule fois, une comparaison
entre un modèle appris et un modèle métier porte réellement sur les modèles, et non sur des
préparations de données divergentes. Cette propriété, invisible dans l'interface, est ce qui rend
la plateforme utilisable comme instrument de recherche et non seulement comme outil de
consultation.

\section{État de mise à disposition}

Les deux dépôts sont publics, sous licence MIT, qui n'impose ni réciprocité ni autorisation
préalable à la reprise :

\begin{itemize}
  \item l'entrepôt, \url{https://github.com/xairon/hubeau_data_integration} ;
  \item la plateforme, \url{https://github.com/xairon/time-serie-explo}.
\end{itemize}

\noindent
Chacun porte sa documentation d'installation, éprouvée depuis une machine vierge, et son manuel
d'exploitation. Reprendre ces travaux ne suppose donc ni demande d'accès, ni contact préalable
avec leur auteur.

La plateforme est déployée (chapitre~\ref{ch:observatoire}), mais son usage demeure à ce jour
circonscrit au cercle du projet. L'accès est nominatif : il n'existe ni inscription libre ni
authentification déléguée, les comptes étant créés par un administrateur. Ce choix, adapté à un
outil scientifique manipulant une chaîne de traitement en évolution, constitue le principal frein
à une diffusion plus large, et le premier levier à actionner pour l'obtenir.

Le socle a par ailleurs été présenté aux équipes du projet ANR \textsc{Aida}, dont le domaine
d'expérimentation recouvre en partie celui traité ici, afin que celles qui en auraient l'usage
n'aient pas à le reconstruire.

Plusieurs pistes ouvertes au cours de ces échanges n'ont pas été suivies. Le rapprochement avec
les modèles à réservoirs du BRGM en est la principale : elle donnerait une seconde référence de
comparaison, indépendante de celle construite ici. Quant à l'élargissement des usages, il tient à
une politique d'accès et à un accompagnement, non à des développements supplémentaires.
